\clearpage
\begin{explanationbox}
\textbf{先看实数模型：椭圆 $x^2/8+y^2/4=1$，直线 $x+y-1=0$。}
\begin{center}\begin{tikzpicture}[font=\small]
\draw[black!30,rounded corners=2pt] (0,0) rectangle (14.4,8.2);
\node[anchor=west,font=\small\bfseries] at (.5,7.72) {消元不是另造点：两个实根分别对应原来的两个交点};
% MODEL 8 4 1 1 -1
\begin{scope}[shift={(7.2,3.65)},scale=1.08]
\coordinate (P) at (-0.8968052533,1.8968052533);
\coordinate (Q) at (2.2301385866,-1.2301385866);
\coordinate (M) at (0.6666666667,0.3333333333);
\coordinate (O) at (0.0000000000,0.0000000000);
\draw[blue,very thick] (0,0) ellipse (2.8284271247 and 2.0000000000);
\draw[-Stealth,black!45] (-3.3,0)--(3.3,0) node[right] {$x$};\draw[-Stealth,black!45] (0,-2.3)--(0,2.3) node[above] {$y$};
\draw[orange!85!black,very thick] (P)--(Q);\draw[green!45!black,dashed] (P)--(P|-O);\draw[green!45!black,dashed] (Q)--(Q|-O);\fill[black] (P) circle (1.8pt) node[above left=5pt] {$P(x_1,y_1)$};
\fill[black] (Q) circle (1.8pt) node[below right=5pt] {$Q(x_2,y_2)$};
\node[below=7pt,text=green!45!black] at (P|-O) {$x_1$};\node[above=7pt,text=green!45!black] at (Q|-O) {$x_2$};
\end{scope}
\node[anchor=west] at (.8,6.7) {$x^2/8+y^2/4=1,\quad x+y-1=0$};\node[anchor=west] at (.8,.6) {两交点的横坐标，正是消元方程的两个实根};
\end{tikzpicture}\end{center}

将 $y=1-x$ 代入，约去公因数后得到 $3x^2-4x-6=0$。
所以横坐标虽未求出，却已有 $x_1+x_2=4/3$、$x_1x_2=-2$。
每一个实根 $x_i$ 都能\textbf{唯一还原} $y_i=1-x_i$；韦达公式研究的就是这两个交点的坐标。
若只要求坐标的和与积、中点或弦长，通常无需把两个根分别写出来。

\tcbbreak
\textbf{先确认交点存在。固定椭圆，只上下移动水平直线。}
\begin{center}\begin{tikzpicture}[font=\small]
\draw[black!30,rounded corners=2pt] (0,0) rectangle (14.4,8.2);
\node[anchor=west,font=\small\bfseries] at (.5,7.72) {先有实交点，再谈弦：割线、切线、无交点三种状态};
% MODEL 9 4 0 1 -1
\begin{scope}[shift={(7.0,3.2)},scale=1.1]
\draw[blue,very thick] (0,0) ellipse (3.0000000000 and 2.0000000000);
\draw[-Stealth,black!45] (-3.5,0)--(3.5,0) node[right] {$x$};\draw[-Stealth,black!45] (0,-1.9)--(0,2.9) node[above] {$y$};
\draw[orange!85!black,thick,solid] (-3.4,1)--(3.4,1);
\draw[purple,thick,dashed] (-3.4,2)--(3.4,2);
\draw[red!75!black,thick,densely dotted] (-3.4,2.6)--(3.4,2.6);
\coordinate (P) at (-2.5980762114,1.0000000000);
\coordinate (Q) at (2.5980762114,1.0000000000);
\coordinate (T) at (0.0000000000,2.0000000000);
\fill (P) circle (2pt);\fill (Q) circle (2pt);\fill (T) circle (2pt);
\end{scope}
\node[anchor=west,text=orange!85!black] at (11.0,4.3) {两个交点};
\node[anchor=west,text=purple] at (11.0,5.4) {一个切点};
\node[anchor=west,text=red!75!black] at (11.0,6.06) {没有交点};
\node[anchor=west] at (.8,.6) {固定椭圆 $x^2/9+y^2/4=1$，只改变水平直线的高度};
\end{tikzpicture}\end{center}

本图 $\alpha=9,\beta=4$，直线为 $y=c$，故 $D=4,H=4-c^2$。
$c=1,2,13/5$ 分别对应 $H>0,H=0,H<0$。
没有实交点时，代数上的根与系数恒等关系不能当作实弦的几何数据。

\tcbbreak
\textbf{为什么要检查消元方向？竖直线不能用横坐标区分交点。}
\begin{center}\begin{tikzpicture}[font=\small]
\draw[black!30,rounded corners=2pt] (0,0) rectangle (14.4,8.2);
\node[anchor=west,font=\small\bfseries] at (.5,7.72) {竖直弦并未消失：横坐标根差为零，不代表弦长为零};
% MODEL 9 4 1 0 -1
\begin{scope}[shift={(7.2,3.65)},scale=1.08]
\draw[blue,very thick] (0,0) ellipse (3.0000000000 and 2.0000000000);
\draw[-Stealth,black!45] (-3.4,0)--(3.4,0) node[right] {$x$};\draw[-Stealth,black!45] (0,-2.3)--(0,2.3) node[above] {$y$};
\coordinate (P) at (1.0000000000,1.8856180832);
\coordinate (Q) at (1.0000000000,-1.8856180832);
\draw[orange!85!black,very thick] (1,-2.3)--(1,2.3);\fill[black] (P) circle (1.8pt) node[above right=6pt] {$P$};
\fill[black] (Q) circle (1.8pt) node[below right=6pt] {$Q$};

\end{scope}
\node[anchor=west] at (.8,6.7) {$x_1=x_2=1$，但 $P\ne Q$};\node[anchor=west] at (.8,.6) {这里 $B=0$，不能用关于 $x$ 的消元判别式判断相切};
\end{tikzpicture}\end{center}

图中 $x=1$ 与椭圆 $x^2/9+y^2/4=1$ 有两个交点，均满足 $x_1=x_2=1$。
此时 $B=0$，不能解出 $y=-(Ax+C)/B$；机械写出的 $x$ 多项式为 $9(x-1)^2=0$，
其“重根”\textbf{不是相切证据}，因为它已不能与原方程组等价对应。
正确做法是固定 $x=1$，解 $y^2=32/9$；或者直接用 $H=9-1=8>0$ 判断。
\end{explanationbox}
